\section{The smallest odd-order finite-gauge separator}\label{sec:sep}

Let \(K\) be finite and \(\omega\in H^3(BK;U(1))\).
For a connected closed oriented three-manifold \(M\), our finite-gauge
normalization is
\begin{equation}\label{eq:dw}
 \mathcal D_{K,\omega}(M):=\frac{Z_{K,\omega}(M)}{Z_{K,\omega}(S^3)}
 =\sum_{\rho\in\Hom(\pi_1 M,K)}
      \big\langle (B\rho)^*\omega,[M]\big\rangle .
\end{equation}
Here \(B\rho\) denotes the induced classifying map from \(M\).
The physical sum uses weight \(1/|K|\) before vacuum division
\cite{DW,FQ}; equivalently sum once per conjugacy orbit with
reciprocal centralizer weight. No division by \(K\) is made in the
right side of \eqref{eq:dw}. Every weight in the results below is one;
in particular, the two computed physical amplitudes are positive.

\subsection{Why no odd nilpotent group can detect the pair}

\begin{theorem}\label{thm:nilp}
Suppose \(p_v=2\), \(\det A_\epsilon\ne0\) and \(K\) is finite
nilpotent of odd order, coprime to \(\det A_\epsilon\).
Let \(\pi:\pi_1 M_{G,\epsilon}\twoheadrightarrow F_{b_G}\) be
the natural graph quotient. Every \(K\)-representation factors
uniquely through \(\pi\). Hence for all \(\omega\),
\[
 Z_{K,\omega}(M_{G,\epsilon})=|K|^{b_G-1},\qquad
 \mathcal D_{K,\omega}(M_{G,\epsilon})=|K|^{b_G}.
\]
The pro-\(\ell\) completion is free pro-\(\ell\) of rank \(b_G\)
for every odd prime \(\ell\nmid\det A_\epsilon\).
\end{theorem}
\begin{proof}
Choose a maximal spanning tree and based port loops with the
two cap loops consecutive in the boundary product. The graph
presentation has a fibre \(h_v\) and cap sections \(a_v,b_v\)
at each vertex, with \(a_v^2=h_v^{-1}\), \(b_v^2=h_v\).
They are the only generators in addition to the ordinary seam
sections and the \(b_G\) stable letters. In a finite group of odd
exponent \(d\) squaring is bijective, with inverse the power
\(u\) for \(2u\equiv1\bmod d\). Consequently the images of \(a_v,b_v\)
are \(h_v^{-u},h_v^u\); they multiply to one and all the required
commutations with \(h_v\) hold in every such image.
The surviving vertex equation is an ordered product of its seam
section images equal to one. Each such section is a conjugate of
\(h_w^{\epsilon_{vw}}\), by the identity on tree seams and by a
stable letter or its inverse otherwise.

Let \(\gamma_1K=K\), \(\gamma_{j+1}K=[K,\gamma_jK]\).
If all fibre images lie in \(\gamma_jK\), conjugations have no
effect in \(\gamma_jK/\gamma_{j+1}K\). In this finite abelian
quotient the vertex equations are
\(A_\epsilon\overline{\mathbf h}=0\).
The adjugate annihilates \(\overline{\mathbf h}\) by
\(\det A_\epsilon\), which is invertible on the quotient by
coprimality. Thus the fibres lie in \(\gamma_{j+1}K\). Beginning
at \(j=1\), nilpotence forces all fibres, then all vertex groups,
to have trivial image. The remaining representation is arbitrary
on the free stable letters.

The projecting map is realized topologically by collapsing the
capped blocks and using intervals on the seam collars. Thus the
classifying maps in this calculation factor through a graph,
where degree-three cohomology vanishes. Formula \eqref{eq:dw}
now gives the stated numbers. The same argument applies to all
finite \(\ell\)-groups for the indicated primes, and identifies
the corresponding pro-\(\ell\) completions by their universal
property.
\end{proof}

For the pair \(M_0,M_1\), the determinant moduli of the signed
adjacencies are \(2^6\) and \(2^8\). Thus all finite odd
nilpotent theories, including all their twists, have the
common vacuum-divided value \(|K|^9\). This conclusion concerns
a whole finite-group class, not just abelianization.

\subsection{A conjugation channel at order twenty-one}

\begin{theorem}\label{thm:sep}
Let \(K=\FF_7\rtimes C_3\), with \(C_3\) acting on the additive
\(\FF_7\) by multiplication by \(2\). For every
\(\omega\in H^3(BK;U(1))\),
\[
 \mathcal D_{K,\omega}(M_0)=2{,}405{,}357{,}452{,}449,\qquad
 \mathcal D_{K,\omega}(M_1)=2{,}727{,}379{,}236{,}309 .
\]
There are three twists and each gives these same two values.
No odd-order finite gauge group of smaller order distinguishes
this pair, with any twist.
\end{theorem}
\begin{proof}
Write the multiplication as
\((x,e)(y,f)=(x+2^e y,e+f)\), with \(e,f\in\ZZ/3\).
Under the projection to \(C_3\) every representation kills all
vertex groups by Theorem~\ref{thm:nilp}; in particular all fibre
images lie in \(\FF_7\). The argument only uses invertibility of
the signed adjacency modulo three for this abelian quotient.
The cap images are then uniquely determined by their fibres,
and both the cap and seam sections lie in the abelian normal
subgroup. A stable letter \((u,e)\) conjugates a fibre by \(2^e\),
independently of \(u\).

Index each part of \(K_{4,4}\) by \(0,1,2,3\). Use the spanning
tree consisting of row zero and column zero. Tree exponents
are zero; the other nine exponents \(e_{ij}\) are independently
in \(\ZZ/3\). With \(b_{ij}=B(D)_{ij}\in\{+1,-1\}\), the fibre
relations, written additively in \(\FF_7\), are
\begin{equation}\label{eq:rank}
  A^+(\mathbf e)\,h_B=0,\qquad
  A^-(\mathbf e)^{\mathsf T} h_A=0,\quad
  A^+_{ij}=b_{ij}2^{e_{ij}},\quad
  A^-_{ij}=b_{ij}2^{-e_{ij}} .
\end{equation}
The reverse power comes from the inverse stable letter at the
other end of the same seam. The signs are the exponents
\(\epsilon_{ij}\) on the adjacent fibre, not new choices of
stable letters. After specifying both fibre vectors, all vertex
relations and every seam relation hold: the defined cap roots
commute with their fibres and the vertex products in the normal
subgroup are precisely the two systems in \eqref{eq:rank}.
There are no additional stable-letter restrictions, by the
graph-of-groups presentation. Hence the number of representations
at fixed exponents is
\(7^9\,7^\nu\), where
\(\nu=8-\rank_{\FF_7} A^+-\rank_{\FF_7}A^-\).
The factor \(7^9\) is from the independent normal components
of the nine stable letters. These are raw based homomorphism
counts, not counts modulo conjugation.

Here is the complete finite enumeration. Both first rows and
columns equal one, so subtracting their pivot reduces each
rank to one plus the rank of the \(3\times3\) matrix with entries
\(b_{ij}2^{\pm e_{ij}}-1\) for \(i,j\ge1\). For a \(3\times3\)
matrix the rank is \(3\) if its determinant is nonzero, else
\(2\) if a \(2\times2\) minor survives, else \(1\) if an entry
survives, else \(0\). Enumerate all nine exponents by the Cartesian
product of three choices, in row-major order. The following
histogram lists all nonempty bins:
\begin{center}
\begin{tabular}{lrrrr}
\toprule
seam matrix & \(\nu=0\) & \(\nu=1\) & \(\nu=2\) &
          \(\sum_{\mathbf e}7^\nu\)\\
\midrule
\(B(I_3)\)       & 14415 & 5070 & 198 & 59607\\
\(B(J_3-I_3)\)   & 13687 & 5712 & 284 & 67587\\
\bottomrule
\end{tabular}
\end{center}
For explicit reproducibility the entire arithmetic in each loop
is
\[
 (u^+,u^-)=(1,1),(2,4),(4,2)\quad
 \longmapsto\quad
 (b_{ij}u^+-1,b_{ij}u^--1)\pmod 7
\]
entry by entry. Start each histogram at zero; for each array of
these nine choices increment index
\(6-\rank(M^+)-\rank(M^-)\) once.
The accompanying script additionally verifies every rank by
ordinary modular elimination on the original \(4\times4\)
matrices. Each row sums to \(3^9=19683\), and the last column is
the respective dot product with \((1,7,49)\). Multiplication by
\(7^9\) gives the displayed exact vacuum-divided values in the
untwisted case. This finite enumeration is part of the calculation,
not an asymptotic experiment.

To address twists, consider \(1\to C_7\to K\to C_3\to1\).
For \(q>0\), \(H^q(C_7;U(1))\) is \(7\)-primary, so its
positive-degree \(C_3\) cohomology is zero by averaging.
The \(q=3\) group is \(\ZZ/7\) and the automorphism acts as
\(2^2=4\) (or its inverse for the opposite pullback convention).
Indeed \(H^3(C_7;U(1))\cong H^4(C_7;\ZZ)\) is generated by
the square of the degree-two cyclic class. There are no invariants.
Thus in total degree three the Lyndon--Hochschild--Serre terms
leave only the inflation from \(H^3(C_3;U(1))\cong\ZZ/3\).
The section \(C_3\hookrightarrow K\) makes inflation injective,
so it supplies all three twists. Each projected \(C_3\)
character factors through the backbone; the pullback of its
degree-three class therefore vanishes on the manifold.
All weights in \eqref{eq:dw} are exactly one for every twist.

Finally the positive odd integers below \(21\) are \(1\),
the primes below \(21\), \(9\) and \(15\). Groups of these
orders are nilpotent: order \(9\) is a prime power; at order
\(15\) both Sylow subgroups are normal by their counting
congruences, so the group is cyclic. Theorem~\ref{thm:nilp}
therefore applies to every smaller odd gauge group for this
pair, proving minimality.
\end{proof}

The mechanism in \eqref{eq:rank} is a fixed normal-module
character that survives conjugation by graph letters.
The lower central filtration of an odd nilpotent target instead
kills it. The result is not a classification of all nonnilpotent
groups or of detectors across other seam signings.
